\documentclass[12pt, a4paper]{article}

% Paquetes requeridos para el circuito
\usepackage[utf8]{inputenc}
\usepackage[T1]{fontenc}
\usepackage{circuitikz} 

\begin{document}

\begin{center}
    \begin{circuitikz}[american, scale=1.0, transform shape]
        
        % 1. Raíl inferior común (nodo 'b')
        \draw (0,0) -- (8.5,0) node[circ] {} node[right] {b};
        
        % 2. Fuente independiente de corriente I entre 'b' y 'c'
        % Añadimos i=$i_1$ para indicar explícitamente que la corriente de esta rama es i1
        \draw (0,0) node[circ] {} to[I, l=$I$, i=$i_1$] (0,3) node[circ] {} node[above] {c};
        
        % 3. Resistencia R en paralelo (entre 'c' y 'b')
        \draw (0,3) -- (2.5,3) to[R, l=$R$] (2.5,0) node[circ] {};
        
        % 4. Fuente dependiente de tensión A*i1 entre 'c' y 'd'
        % Utilizamos 'cV' para la fuente dependiente (rombo)
        \draw (2.5,3) to[cV, l=$A i_1$] (5.5,3) node[circ] {} node[above] {d};
        
        % 5. Resistencia R entre 'd' y 'a' (Resistencia de salida del modelo)
        \draw (5.5,3) to[R, l=$R$] (8.5,3) node[circ] {} node[above] {a};
        
        % 6. Potenciómetro RL entre 'a' y 'b' como carga (usamos vR para resistencia variable)
        \draw (8.5,3) to[vR, l=$R_L$] (8.5,0) node[circ] {};
        
    \end{circuitikz}
\end{center}

\end{document}